\chapter{Dictionary}
\newpage
\noindent\textbf{Airfoil} The shaped cross-section of a wing. It makes the air move differently over the top and bottom of the wing and creates lift.

\vspace{0.45em}
\noindent\textbf{Angle of attack} The angle at which the wing meets the air flow. A greater angle usually creates more lift, but also more drag.

\vspace{0.45em}
\noindent\textbf{Centre of rotation} The point on the boomerang around which the boomerang spins. It belongs to the boomerang itself and stays fixed on the boomerang while it rotates.

\vspace{0.45em}
\noindent\textbf{Construction disk} A drawing tool used to understand where different parts of the boomerang meet faster or slower air flow. It helps you compare how shape, airfoil, dihedral, weight or carving may change the flight.

\vspace{0.45em}
\noindent\textbf{Dihedral} A bend or tilt in a wing. Positive and negative dihedral can change the amount of lay-over and the height of the flight.

\vspace{0.45em}
\noindent\textbf{Eagle catch} A catch imitating a hunting eagle performed with a multi-winged boomerang. You catch the boomerang with one hand from above by closing your fingers around the centre of the boomerang.

\vspace{0.45em}
\noindent\textbf{Foot catch} Just what it sounds like. A catch performed by catching the boomerang with your feet. Your feet are not allowed to touch the ground while catching, and the boomerang has to be held for a few seconds.

\vspace{0.45em}
\noindent\textbf{Forward speed ($V_t$)} The speed with which the boomerang moves forward through the air.

\vspace{0.45em}
\noindent\textbf{Gyroscopic precession} The delayed response of a spinning object to a tilting force. On a boomerang, a lift force has its main effect about 90° later in the direction of rotation.

\vspace{0.45em}
\noindent\textbf{Instant centre} The point on the construction disk that represents the momentary centre of the combined forward and spinning motion. When forward speed is high, the boomerang's centre of rotation is placed farther away from the instant centre.

\vspace{0.45em}
\noindent\textbf{Lay-over} Lay-over generally refers to the boomerang's angle to vertical. The greater the angle, the more lay-over. It can also refer to the act: the boomerang lays over from a vertical position to a horizontal position.

\vspace{0.45em}
\noindent\textbf{Lift} The aerodynamic force created by the airfoils. More air speed and a better angle to the air flow usually create more lift.

\vspace{0.45em}
\noindent\textbf{Rotation speed ($V_r$)} The speed caused by the boomerang spinning around its centre of rotation.

\vspace{0.45em}
\noindent\textbf{Tunnel catch} A catch performed by standing with both feet on ground, while catching the boomerang between the legs. This means you have to have one hand behind your legs and one hand in front and then ``sandwich'' the boomerang while passing it between your legs with one hand. 
